\documentclass[11pt]{article}
\usepackage[margin=1in]{geometry}
\usepackage{amsmath,amssymb,amsthm}
\usepackage{hyperref}
\usepackage[T1]{fontenc}
\usepackage[utf8]{inputenc}
\usepackage{enumitem}
\usepackage{array}
\usepackage{longtable}

\newtheorem{definition}{Definition}
\newtheorem{designprinciple}{Design Principle}
\newtheorem{proposition}{Proposition}
\newtheorem{conjecture}{Conjecture}
\newtheorem{example}{Example}

\title{From ***plain Specifications to MeTTa/PeTTa and Rholang:\
A Typed Atomspace IR Proposal}
\author{OmegaClaw / ZeroBot working notes for Ben Goertzel}
\date{2026-06-29}

\begin{document}
\maketitle

\begin{abstract}
This note reviews a short design discussion about whether software
specifications written in the ***plain style can be mapped into MeTTa, initially
PeTTa as a MeTTa dialect, and later into Rholang and MeTTa-IL.  The conclusion
is positive but conditional: the viable target is not a magic compiler from
arbitrary English to correct code, but an assisted semantic lowering pipeline
from structured specifications to an inspectable typed logical intermediate
representation.  The preferred intermediate representation is an
Atomspace-style typed logical graph, designed as a logical mirror of MeTTa and
constrained by Curry-Howard and OSLF-style correspondences.  This IR should make
MeTTa/PeTTa generation relatively natural, keep Rholang generation grounded in
explicit process semantics, and open the door to later PLN-based analysis of
requirements, contradictions, underspecification, test coverage, and confidence.
\end{abstract}

\section{Provenance and purpose}

This document summarizes and expands a 2026-06-29 discussion between Ben
Goertzel, ProtomegaTron/OmegaClaw, and ZeroBot/OpenClaw about using
\texttt{***plain} software specifications as input to a semantic toolchain.
The motivating links were the ***plain documentation at
\url{https://plainlang.org/docs/} and the Plain GitHub organization at
\url{https://github.com/plainlang}.

The immediate question was whether one could develop a tool that maps plain
software specifications into MeTTa, using PeTTa as an initial target dialect,
and then into Rholang as well.  A later extension would target MeTTa-IL.  Ben
then suggested that an Atomspace-based logic representation might be the right
intermediate representation, for two reasons:
\begin{enumerate}[label=(\arabic*)]
  \item if the logic representation is a good logical mirror of MeTTa, in the
        sense suggested by Curry-Howard and Meredith/Stay OSLF-style type
        mappings, then mapping the IR to MeTTa code becomes natural rather than
        ad hoc;
  \item such an IR opens the door to later PLN-based analysis of specifications.
\end{enumerate}

The purpose of this note is to turn that conversational design intuition into a
more detailed information package for collaborators who may suggest designs,
identify risks, or help build a prototype.

\section{What ***plain contributes}

The public ***plain documentation presents ***plain as a specification language
for writing software requirements in a clear structured format.  A file has
optional YAML frontmatter followed by standardized sections.  The core sections
include:
\begin{itemize}
  \item \texttt{***definitions***}, declaring named concepts such as
        \texttt{:App:}, \texttt{:User:}, \texttt{:Task:}, and their attributes
        or constraints;
  \item \texttt{***implementation reqs***}, giving non-functional and
        implementation-oriented constraints such as language, framework,
        architecture, naming conventions, and environment;
  \item \texttt{***test reqs***}, specifying how conformance tests should be
        structured and executed;
  \item \texttt{***functional specs***}, listing small, testable functional
        requirements; and
  \item nested \texttt{***acceptance tests***}, giving observable criteria under
        individual functional requirements.
\end{itemize}

Important properties for this proposal are:
\begin{enumerate}
  \item concepts are explicitly named using a lightweight notation such as
        \texttt{:Task:};
  \item definitions can refer to other definitions;
  \item attributes and constraints can be represented as nested bullets;
  \item functional requirements are already encouraged to be small and
        testable;
  \item acceptance tests are explicitly linked to requirements; and
  \item imported modules and template libraries can contribute reusable
        definitions and constraints.
\end{enumerate}

These features make ***plain a promising human-readable front end for semantic
software specification.  The crucial limitation is that many bodies are still
natural language.  Therefore, the MVP should avoid pretending that arbitrary
English has been fully formalized.

\section{Main viability claim}

\begin{proposition}[Viability as assisted semantic lowering]
A toolchain from ***plain to MeTTa/PeTTa and Rholang is viable if it is framed as
an assisted semantic lowering pipeline with an inspectable typed logical IR, not
as an automatic compiler from arbitrary natural language into correct executable
systems.
\end{proposition}

The proposed pipeline is:
\begin{verbatim}
.plain source
  -> parsed Plain AST
  -> typed Atomspace / logical Spec IR
  -> MeTTa / PeTTa emitter
  -> PLN and consistency-analysis passes
  -> Rholang process-model emitter
  -> later MeTTa-IL emitter
\end{verbatim}

The MVP can deliberately keep natural-language interpretation shallow.  It can
parse explicit section structure, concept names, bullets, nesting, provenance,
and some constrained patterns, while preserving ambiguous text as quoted claims
or obligations for later human/LLM refinement.

\section{Why an Atomspace-style logical IR is attractive}

\begin{designprinciple}[Typed logical mirror]
The central IR should be a typed logical graph that mirrors the relevant logical
and type-theoretic structure of MeTTa closely enough that MeTTa emission is a
principled rendering, not a backend-specific rewrite trick.
\end{designprinciple}

An Atomspace-style IR is attractive because it can represent concepts,
relations, propositions, procedures, tests, and truth/confidence annotations in a
single hypergraph-like substrate.  It can also preserve provenance and partial
formalization status without forcing premature executable precision.

\begin{definition}[Spec atom]
A spec atom is a typed IR element
\[
  a = \langle id, kind, type, args, provenance, status, truth \rangle
\]
where \textit{kind} distinguishes concepts, predicates, links, requirements,
processes, tests, and annotations; \textit{type} records the logical or
computational type; \textit{args} are outgoing targets; \textit{provenance}
links back to the source file, section, bullet, and textual span;
\textit{status} records whether the atom is parsed, inferred, reviewed,
rejected, or executable; and \textit{truth} can later store PLN-style strength
and confidence.
\end{definition}

\begin{definition}[Plain-to-IR lowering]
A Plain-to-IR lowering function is a partial translation
\[
  L : PlainAST \to (G, W)
\]
where \(G\) is a typed Atomspace-style graph and \(W\) is a set of warnings,
ambiguities, unresolved references, and human-review obligations.
\end{definition}

The word ``partial'' matters.  The first version should translate what is
structural and typed, preserve what is natural-language and ambiguous, and make
uncertainty visible.

\section{OSLF, Curry-Howard, and the type layer}

Ben's key refinement was that even without full English-to-logic parsing, the
system should get type relationships right so that OSLF-style correspondences
can work later.  This suggests that the type layer is not decorative metadata;
it is the spine of the architecture.

\begin{designprinciple}[Types before prose semantics]
For the MVP, prioritize a correct concept/type/proposition layer over ambitious
natural-language semantic parsing.  Natural language can remain attached as
quoted evidence, but type relationships should be explicit, checkable, and
mapped coherently into MeTTa/PeTTa.
\end{designprinciple}

A useful three-level distinction is:
\begin{enumerate}
  \item \textbf{Concept layer.} Plain concepts such as \texttt{:Task:},
        \texttt{:User:}, \texttt{:TaskList:}, \texttt{:Implementation:}, and
        \texttt{:ConformanceTests:}.
  \item \textbf{Logical type graph.} Atomspace-like nodes and links expressing
        inheritance, instancehood, products, functions, predicates,
        propositions, obligations, and tests.
  \item \textbf{Operational rendering.} MeTTa/PeTTa declarations and rules,
        Rholang processes and channels, and later MeTTa-IL constructs.
\end{enumerate}

The Curry-Howard intuition is that propositions, types, proofs, programs, and
process contracts should be aligned rather than represented in incompatible
layers.  A functional spec can be treated as an obligation/proposition whose
realizer is code and whose evidence includes acceptance-test traces.  A concept
can be represented both as a type and as an object of reasoning.  A test can be a
claim about observable behavior and also a process that attempts to produce
witnesses or counterexamples.

\section{Candidate IR vocabulary}

The following table gives an initial vocabulary.  The exact Atom names are not
final; the important issue is the separation of roles.

\begin{longtable}{|p{1.5in}|p{2.1in}|p{2.5in}|}
\hline
Plain element & IR role & Possible MeTTa/PeTTa rendering \\
\hline
\texttt{:Concept:} definition & Concept/type declaration & \texttt{(: Concept Type)} or domain-specific type atom \\
\hline
Concept reference & Reference edge, possibly typed & \texttt{(refers-to Req Concept)} \\
\hline
Nested attribute & Attribute predicate or dependent field & \texttt{(has-attribute Task Name Required)} \\
\hline
Implementation requirement & Non-functional constraint / realization hint & \texttt{(implementation-constraint Req Python)} \\
\hline
Test requirement & Conformance-test policy & \texttt{(test-policy App Unittest)} \\
\hline
Functional spec bullet & Requirement proposition / obligation & \texttt{(: FS1 Requirement)} and \texttt{(requires FS1 ...)} \\
\hline
Acceptance test & Observable evidence condition & \texttt{(acceptance-test FS1 AT1)} \\
\hline
Import/template & Module dependency and namespace source & \texttt{(imports Spec Module)} \\
\hline
Natural-language residue & Quoted claim needing interpretation & \texttt{(source-text FS1 "...")} \\
\hline
\end{longtable}

An Atomspace-like version might use nodes and links such as:
\begin{verbatim}
ConceptNode("Task")
ConceptNode("User")
InheritanceLink(ConceptNode("Task"), ConceptNode("PlainConcept"))
EvaluationLink(PredicateNode("has-attribute"),
  ListLink(ConceptNode("Task"), ConceptNode("Name"), ConceptNode("Required")))
EvaluationLink(PredicateNode("functional-spec"),
  ListLink(ConceptNode("FS3"), ConceptNode("User"),
           PredicateNode("can-add"), ConceptNode("Task")))
\end{verbatim}

A MeTTa/PeTTa rendering of the same structure could be:
\begin{verbatim}
(: Task PlainConcept)
(: User PlainConcept)
(: TaskList PlainConcept)
(has-attribute Task Name Required)
(functional-spec FS3)
(actor FS3 User)
(action FS3 add)
(object FS3 Task)
(constraint FS3 (valid Task))
(acceptance-test FS3 AT3_1)
(source-text FS3 ":User: should be able to add :Task:...")
\end{verbatim}

\section{Worked miniature example}

A small ***plain fragment might be:
\begin{verbatim}
***definitions***
- :App: is a console task manager application.
- :User: is the user of :App:.
- :Task: describes an activity that needs to be done by :User:.
  - Name - a short required description of :Task:.
- :TaskList: is a list of :Task: items.
  - Initially :TaskList: should be empty.

***implementation reqs***
- :Implementation: should be in Python.

***test reqs***
- :ConformanceTests: of :App: should be implemented using Unittest.

***functional specs***
- :User: should be able to add :Task:. Only valid :Task: items can be added.
  ***acceptance tests***
  - Adding a valid :Task: should make it appear in :TaskList:.
\end{verbatim}

A conservative IR should produce at least:
\begin{enumerate}
  \item a concept/type environment containing \texttt{App}, \texttt{User},
        \texttt{Task}, \texttt{TaskList}, \texttt{Implementation}, and
        \texttt{ConformanceTests};
  \item relationship atoms such as \texttt{user-of(User, App)},
        \texttt{list-of(TaskList, Task)}, \texttt{has-field(Task, Name)};
  \item an initial-state constraint \texttt{empty(TaskList)};
  \item an implementation constraint \texttt{target-language(Python)};
  \item a test policy \texttt{framework(Unittest)};
  \item a requirement atom for add-task behavior; and
  \item an acceptance-test atom linked to the requirement.
\end{enumerate}

A possible PeTTa-oriented representation is:
\begin{verbatim}
(: App PlainConcept)
(: User PlainConcept)
(: Task PlainConcept)
(: TaskList PlainConcept)
(: AddTask Requirement)
(: AddTaskAcceptance AcceptanceTest)

(user-of User App)
(list-of TaskList Task)
(has-field Task Name)
(required-field Task Name)
(initial-state TaskList Empty)
(target-language Implementation Python)
(test-framework ConformanceTests Unittest)

(requirement AddTask)
(actor AddTask User)
(action AddTask add)
(object AddTask Task)
(precondition AddTask (valid Task))
(postcondition AddTask (member Task TaskList))

(acceptance-test-of AddTaskAcceptance AddTask)
(observes AddTaskAcceptance (member Task TaskList))
(source-text AddTaskAcceptance
  "Adding a valid :Task: should make it appear in :TaskList:.")
\end{verbatim}

This is intentionally not yet an executable task manager.  It is a scrutable
semantic skeleton from which code-generation, reasoning, and test-generation
steps can proceed.

\section{Mapping to MeTTa and PeTTa}

The MeTTa/PeTTa path is the most direct because the proposed IR and MeTTa share
a bias toward typed symbolic structures, hypergraph-like representation, and
executable logic.  PeTTa, as a MeTTa dialect or restricted profile for early
experiments, can focus on a disciplined subset:
\begin{itemize}
  \item concept declarations;
  \item type declarations;
  \item relation atoms;
  \item requirement and acceptance-test atoms;
  \item provenance atoms;
  \item simple rule templates for validation and test generation; and
  \item truth-value/confidence slots reserved for later PLN integration.
\end{itemize}

The first emitter should not try to synthesize full algorithms.  It should emit
semantic declarations and obligations.  Examples of early useful outputs:
\begin{enumerate}
  \item undefined-concept warnings;
  \item duplicate-concept warnings;
  \item requirement-to-test coverage maps;
  \item type-consistency checks;
  \item a queryable graph of concepts, attributes, and requirements;
  \item PeTTa rules that identify underspecified requirements;
  \item PeTTa rules that identify acceptance tests whose observable predicate is
        not connected to the parent requirement; and
  \item human-readable roundtrip reports.
\end{enumerate}

\section{Mapping to Rholang}

The Rholang mapping is less direct but still meaningful.  Rholang is a process
calculus language, so its natural semantic targets are not merely concepts and
requirements but channels, messages, processes, state transitions, and
observable traces.

A useful mapping discipline is:
\begin{itemize}
  \item concepts become data schemas, channel payload types, or process roles;
  \item functional specs become process contracts or transition obligations;
  \item user actions become receives/sends on channels;
  \item application state becomes a collection of process-held names or stores;
  \item acceptance tests become trace assertions over observable behavior;
  \item implementation requirements become backend constraints, not semantic
        truths; and
  \item concurrency-relevant requirements become explicit channel/process
        structure.
\end{itemize}

For the task example, a Rholang-oriented IR might identify:
\begin{verbatim}
Role(User)
Process(App)
Channel(addTask)
Channel(taskListOut)
Message(addTask, Task)
State(TaskList)
TraceAssertion(after send(addTask, validTask),
               eventually observe(taskListOut, contains(validTask)))
\end{verbatim}

The Rholang emitter should initially target small examples where the operational
model is obvious.  Generating Rholang from arbitrary business prose would be
premature.  However, if the Atomspace IR distinguishes state, roles, actions,
messages, and observations, then Rholang generation becomes a controlled
projection from a process-enriched semantic graph.

\section{PLN-enabled specification analysis}

The Atomspace IR also enables a later PLN analysis layer.  This is not needed
for the first parser, but it should influence the representation now.

Potential PLN/spec-analysis tasks include:
\begin{enumerate}
  \item infer that a requirement is probably underspecified because it mentions
        an action but no object, state change, or observable acceptance test;
  \item infer likely concept equivalences or near-duplicates across imported
        modules;
  \item detect contradictions between implementation requirements;
  \item estimate confidence that acceptance tests cover a requirement;
  \item propagate confidence from ambiguous natural-language parses into later
        code-generation obligations;
  \item identify requirements with no tests or tests with no requirement;
  \item reason about whether a Rholang trace assertion plausibly realizes a
        MeTTa/PeTTa requirement proposition; and
  \item support interactive clarification by asking the human for the smallest
        missing semantic commitment.
\end{enumerate}

This motivates storing provenance, confidence, and interpretation status from
the beginning.

\section{Architecture proposal}

A practical prototype can be organized as follows:
\begin{verbatim}
plainmetta/
  parser/
    plain_sections.py        # sections, bullets, nesting, spans
    concept_refs.py          # :Concept: references and definitions
  ir/
    atoms.py                 # typed Atomspace-like data model
    typecheck.py             # concept/type/proposition checks
    provenance.py            # source spans and status
  emitters/
    metta.py                 # MeTTa/PeTTa emission
    rholang.py               # small process-calculus examples
    reports.py               # HTML/Markdown/LaTeX inspection reports
  analysis/
    coverage.py              # requirements vs acceptance tests
    ambiguity.py             # unresolved NL obligations
    pln_stub.py              # future truth-value and inference hooks
  examples/
    task_manager.plain
    task_manager.metta
    task_manager.rho
\end{verbatim}

\section{MVP plan}

A good MVP should be useful before solving natural-language understanding.

\subsection*{MVP 0: Structural parser and concept graph}
\begin{itemize}
  \item Parse sections, bullets, indentation, source spans, and YAML frontmatter.
  \item Extract concept definitions and concept references.
  \item Detect undefined, duplicate, and unused concepts.
  \item Emit a JSON or S-expression Plain AST.
\end{itemize}

\subsection*{MVP 1: Typed Atomspace Spec IR}
\begin{itemize}
  \item Lower concept definitions, attributes, implementation reqs, test reqs,
        functional specs, and acceptance tests into typed IR atoms.
  \item Preserve natural-language text as quoted source atoms.
  \item Attach provenance, status, and optional confidence.
  \item Add simple type checks and requirement/test coverage reports.
\end{itemize}

\subsection*{MVP 2: PeTTa/MeTTa emitter}
\begin{itemize}
  \item Emit stable PeTTa/MeTTa declarations for concepts, relations,
        requirements, and tests.
  \item Add simple query/rule examples for missing tests, unresolved concepts,
        and suspiciously complex functional specs.
  \item Roundtrip enough structure to make code review meaningful.
\end{itemize}

\subsection*{MVP 3: Rholang micro-emitter}
\begin{itemize}
  \item Select constrained specs with explicit actors, actions, objects, state,
        and observations.
  \item Generate small Rholang process skeletons and trace assertions.
  \item Treat this as process-model generation, not full application synthesis.
\end{itemize}

\subsection*{MVP 4: PLN-ready analysis hooks}
\begin{itemize}
  \item Add truth-value slots and confidence propagation conventions.
  \item Represent ambiguous interpretations as competing atoms with provenance.
  \item Prototype one or two PLN-style checks, such as likely missing acceptance
        tests or contradictory implementation constraints.
\end{itemize}

\section{Key research and engineering questions}

For collaborators, the most valuable questions are:
\begin{enumerate}
  \item What is the smallest typed IR that is genuinely a logical mirror of the
        relevant MeTTa fragment?
  \item How should OSLF-style type mappings be represented concretely enough to
        guide emission and checking?
  \item Which Atomspace links and node categories should be treated as primitive
        in the Spec IR, and which should be macros?
  \item Should PeTTa be a narrow profile of MeTTa syntax, a library convention,
        or a more explicit dialect with its own schema?
  \item What restricted natural-language patterns are worth supporting in the
        MVP, and which should remain quoted obligations?
  \item What is the correct semantic relation between a requirement proposition,
        a program realizer, and an acceptance-test trace?
  \item How much process structure must be present before Rholang generation is
        meaningful?
  \item How should uncertainty and provenance be encoded so that PLN can later
        reason over the graph without confusing source claims, inferred claims,
        decisions, and executable artifacts?
\end{enumerate}

\section{Risks}

\begin{enumerate}
  \item \textbf{Overclaiming automation.} The project should not be marketed as
        automatic formalization of arbitrary English.  It is an assisted
        formalization and lowering tool.
  \item \textbf{Premature backend commitment.} If the IR is too close to one
        backend, Rholang and MeTTa-IL may become awkward.  The IR should be
        logical first, backend-renderable second.
  \item \textbf{Weak type discipline.} If concept/type/proposition distinctions
        are fuzzy, the later OSLF/Curry-Howard story will become cosmetic rather
        than structural.
  \item \textbf{Natural-language residue.} Quoted text must not be silently
        treated as a verified logical formula.  Status and confidence must be
        explicit.
  \item \textbf{Rholang mismatch.} Rholang generation requires operational
        process semantics.  Not every requirement has enough process structure.
  \item \textbf{PLN pollution.} PLN should reason over typed, provenance-aware
        claims, not over undifferentiated strings.
\end{enumerate}

\section{Recommended near-term deliverable}

The recommended first deliverable is a small repository containing:
\begin{enumerate}
  \item a parser for a useful subset of ***plain;
  \item a typed Atomspace-style Spec IR schema;
  \item one task-manager example from ***plain to IR to PeTTa/MeTTa;
  \item one small actor/process example from IR to Rholang skeleton;
  \item an inspection report showing concept graph, requirement coverage,
        acceptance-test links, warnings, and unresolved natural-language
        obligations; and
  \item a short design note explaining how the type layer is intended to support
        OSLF/Curry-Howard correspondences and later PLN analysis.
\end{enumerate}

This deliverable would be small enough to build quickly but rich enough to test
the central architectural bet: that Plain's readable specification structure can
be lowered into an Atomspace-style typed logical graph that naturally emits
MeTTa/PeTTa and can later support Rholang, MeTTa-IL, and PLN analysis.

\section{Bottom line}

The strongest version of the idea is:
\begin{verbatim}
***plain
  -> human-readable structured spec
  -> typed Atomspace logical mirror
  -> PeTTa/MeTTa semantic rendering
  -> PLN-style spec analysis
  -> Rholang process realization for process-suitable fragments
  -> later MeTTa-IL lowering
\end{verbatim}

This keeps the simple path simple: parse concepts, requirements, tests, and
provenance.  But it also avoids a dead-end architecture by making the IR rich
enough for type-theoretic interpretation, MeTTa alignment, PLN analysis, and
process-calculus generation.

\end{document}
